\begin{center}
\textsc{Gauss} to \textsc{Gerling}.\quad Göttingen, 3~October 1823.
\end{center}

\par\noindent\hspace*{\parindent}\mbox{.\ .\ .\ .\ .\ .}\quad Allow me a few further \mbox{remarks.\quad\mbox{.\ .\ .\ .\ .\ .}}
\par\noindent\hspace*{\parindent}1) When comparing two points, I always start with the one
I most fear might fail.
\par\clearpage
\noindent\makebox[0pt][l]{\fontsize{13}{16}\selectfont 390}\hfill{\small\textbf{CORRESPONDENCE WITH GERLING.}}\hfill\null\par\medskip
\par\noindent\hspace*{\parindent}2) If the second point is a heliotrope light and there is less than the greatest confidence
that it will hold, I take readings more often than usual. When the risk
was extremely slight, I sometimes took more than 10, i.e.\ 15,
or even 20 measurements before reading again. One always puts a great deal
at risk in doing this, however. If there is considerable fear that the second
point may fail, I may take a reading after the third or even the second observation, or with
every observation. The same applies if the first point is absent for rather a long
time.
\par\noindent\hspace*{\parindent}3) If I am in the middle of a series of observations in which the first point
has been bisected, I like to have someone look through a hand-held telescope before I release
the alidade to see whether the second point is visible. But if it fails
after the alidade has already been released and does not reappear, the
measurement is not lost for that reason. If need be, I then bisect a third point
and so obtain mixed angles. On my visit to the Brocken this time,
however, this happened only once, when I took 2 times
(Ilefeld-Hohehagen) \(+\) 2 times (Ilefeld-Hils). Such mixed angles are used in the
system just as well as pure ones and wholly pro rata. In 1821 I had
many mixed angles, less often thereafter. At most stations, therefore, I have
auxiliary directions in addition to the principal directions, particularly in
1821, when I had only one heliotrope. I like to choose pointed towers for
auxiliary directions, ones that are reasonably close and, if possible,
reasonably close to the horizontal plane. This last condition could of course
not be met on the Brocken. Huyseburg lies \(1^{0}\,37^{\prime}\), Hüttenrode \(1^{0}\,54^{\prime}\), and a pole
on the Wurmberg that I had erected in 1821 and replaced in 1823
(though at a somewhat different spot) \(2^{0}\,13^{\prime}\) below the horizon.
\par
